\documentclass[10pt,a4paper]{amsart}
\usepackage[margin=19mm]{geometry}
\usepackage{amsmath,amssymb,mathtools}
\usepackage[scr=rsfs,cal=euler]{mathalfa}
\usepackage{xfrac}
\usepackage[unicode,hidelinks,bookmarks=false]{hyperref}
\newcommand{\ie}{i.e.\ }
\newcommand{\cf}{cf.\ }
\newcommand{\resp}{respectively\ }
\newcommand{\Subsection}{\subsection}
\providecommand{\nobreakdash}{\nobreak\hbox{-}}
\setlength{\emergencystretch}{2em}
\multlinegap0pt
\usepackage{bm}
\usepackage{bbm}
\usepackage{dsfont}
\usepackage{xspace}
\usepackage{cancel}
\usepackage{mathtools}
\usepackage{amsthm}
\makeatletter
\@ifundefined{theorem}{\newtheorem{theorem}{Theorem}}{}
\@ifundefined{claim}{\newtheorem{claim}[theorem]{Claim}}{}
\@ifundefined{definition}{\newtheorem{definition}[theorem]{Definition}}{}
\@ifundefined{lemma}{\newtheorem{lemma}[theorem]{Lemma}}{}
\makeatother
\newcommand{\calF}{\mathcal{F}}
\newcommand{\calH}{\mathcal{H}}
\newcommand{\calV}{\mathcal{V}}
\title[Constructions for positional games and applications to domin]{Constructions for positional games and applications to domination games}
\author{Ali Deniz Bagdas, Dennis Clemens, Fabian Hamann, Yannick Mogge}
\date{Source: arXiv:2509.05089v1 (CC BY 4.0)}
\begin{document}
\maketitle

\noindent\emph{Source and licence.} Constructions for positional games and applications to domination games. Version arXiv:2509.05089v1 (CC BY 4.0), distributed under the Creative Commons Attribution 4.0 International licence, as recorded in the arXiv licence metadata for this exact version and shown on the abstract page https://arxiv.org/abs/2509.05089v1. Full rights notice: licenses/arxiv-2509.05089-CC-BY-4.0.txt.\par\smallskip

\noindent\textit{Editorial scope.} Counted unit: the Lemma whose source label is \texttt{lemma:from.H.to.G} in the article (source0.tex lines 402--410, source label \texttt{lemma:from.H.to.G}), delivered together with the article's own complete original proof of it (source0.tex lines 459--555, one \texttt{proof} environment of 97 lines). No mathematics is written, repaired or re-derived: statement and proof are the article's own text line for line. The proof is detached from the statement in the source: the article states the result at lines 402--410 and prints its proof at lines 459--555, and the proof environment's own header names the statement, so the association is the article's own. Required context retained uncounted: def:transference.construction (lines 387--400); clm:transference.claim1 (lines 419--421); clm:transference.claim3 (lines 445--447). Every definition, display and label of those lines is retained unchanged. Omitted from this package and not redistributed: 1--386, 401--401, 411--418, 422--444, 448--458, 556--1453. Nothing inside a retained excerpt is omitted or altered: every retained body is delivered exactly as the article's lines are written, so no omission receipt of any kind accompanies this package. Citations: the retained excerpts carry no citation command at all, so no bibliography entry of the article is transcribed and no reference is redistributed. Reference adaptation: none was needed and none was made; every reference inside the delivered unit resolves inside it, so no unresolved reference or citation is produced. Printed slips kept literal: no printed word, symbol, hypothesis or number of the counted statement or of its proof was altered. Layout and class: the article's own class is \texttt{amsart} with packages including algorithm2e, graphicx, amsmath, amssymb, amsthm, mathrsfs, mathabx, dsfont, xcolor, tikz, hyperref, fontenc, lmodern, microtype; the wrapper is the lane's pinned \texttt{amsart} editorial wrapper at 10pt on A4, which restates the theorem-like environments the retained text needs and adds the article's own macro definitions that the retained text needs (\texttt{\textbackslash calF}, \texttt{\textbackslash calH}, \texttt{\textbackslash calV}), with extra wrapper packages (bm, bbm, dsfont, xspace, cancel, mathtools). The delivered numbering is the unit's own. Assets: no retained span contains a figure environment, a graphics-inclusion command, a PDF-page inclusion, a table or a TikZ picture; no image file is copied into this package, the package declares no asset dependency and no image file is redistributed with it. Licence: version arXiv:2509.05089v1 of this article is distributed under the Creative Commons Attribution 4.0 International licence, as recorded in the arXiv licence metadata for this exact version and shown on the abstract page https://arxiv.org/abs/2509.05089v1. A full rights notice is delivered at licenses/arxiv-2509.05089-CC-BY-4.0.txt. Characters: every retained excerpt is pure ASCII, so the delivered engine, XeLaTeX with its default Latin Modern text font, reports no missing character.
\begin{definition}\label{def:transference.construction}
Let $\calH = (X,\calF)$ be a hypergraph,
let $a\geq 1$ be an integer.
Let $\calV(\calH) \subset \mathcal{P}(X)$ be the set of all vertex covers of $\calH$, and for each $A\in \calV(\calH)$, 
let $X_A$ be a vertex set of size $4a(|X| + |\calV(\calH)|)$,
such that all the sets $X_A$ and $X$ are pairwise disjoint.
We construct a graph $G=G(\calH,a)$
as follows:
\begin{itemize}
\item its vertex set is 
$V(G)=X\cup (\bigcup_{A\in \calV(\calH)} X_A)$, and
\item its edge set is $E(G)=\{vw:~ (\exists A\in \calV(\calH):~v\in A ~ \land ~ w\in X_A)\} \cup \binom{X}{2}$.
\end{itemize}
\end{definition}

\begin{claim}\label{clm:transference.claim1}
Each $F\in \calF$ is a dominating set of $G$.
\end{claim}

\begin{claim}\label{clm:transference.claim3}
If $D\subset X$ is a dominating set in $G$, then there is some $F\in\mathcal{F}$ such that $F\subset D$.
\end{claim}

\begin{lemma}[Transference Lemma]\label{lemma:from.H.to.G}
Let $\calH =(X,\mathcal{F})$ be a hypergraph, 
let $a,b,s,t\in\mathbb{N}$, and let $G=G(\mathcal{H},a)$ be constructed according to Definition~\ref{def:transference.construction}. 
Then the following holds:
\begin{enumerate}
\item[(M)] $\mathcal{H} \in M_i(a:b,s,t) ~ \Leftrightarrow ~ G \in D_{MB,i}(a:b,s,t)$,
\item[(W)] $\mathcal{H} \in W(s,t) ~ \Leftrightarrow ~ G \in D_{WC}(s,t)$.
\end{enumerate}
\end{lemma}

\begin{proof}[Proof of Lemma~\ref{lemma:from.H.to.G}]

\textbf{(M)} 
"$\Rightarrow$" 
Assume that $\mathcal{H} \in M_i(a:b,s,t)$, i.e.~Maker (being the $i^{\text{th}}$ player) has a strategy $\mathcal{S}_M$ to claim a winning set of size at most $s$ within at most $t$ rounds in the $(a:b)$ Maker-Breaker game on $\mathcal{H}$.
Then in the $(a:b)$ Maker-Breaker domination game on $G$, Dominator (being the $i^{\text{th}}$ player) can restrict her moves to $G[X]$ and play according to the strategy $\mathcal{S}_M$. 
This way, Dominator claims a set $F\in \mathcal{F}$ of size at most $s$ within at most $t$ rounds, and thus by Claim~\ref{clm:transference.claim1}, we have $G \in D_{MB,i}(a:b,s,t)$.



"$\Leftarrow$" We do a proof by contraposition. Assume that 
$\mathcal{H} \notin M_i(a:b,s,t)$, and therefore
Breaker has a strategy $\mathcal{S}_B$ 
in the $(a:b)$ Maker-Breaker game on $\calH$ 
(where Maker is the $i^{\text{th}}$ player)
which prevents Maker from claiming a winning set of size
at most $s$ within the first $t$ rounds. Then, in the $(a:b)$ Maker-Breaker domination game on $G$
(where Dominator is the $i^{\text{th}}$ player),
Staller can apply the following strategy:
\begin{itemize}
\item[] \textbf{Stage I:} As long as before Staller's move there is at least one free vertex in $X$,
Staller plays on $G[X]$ only, and he plays according to Breaker's 
strategy $\mathcal{S}_B$. (If Dominator claims an element
outside of $X$, Staller simply pretends that Dominator made an arbitrary move within $X$.)
\item[] \textbf{Stage II:} When Staller enters Stage II, all elements of $X$ are claimed. For each $A\in \calV(\calH)$,
let $X_A'\subset X_A$ be the vertices which are still free at the end of Stage I. Then from now on, for each vertex that Staller can claim,
Staller chooses a vertex from an arbitrary set $X_A'$ in which he still has no vertex claimed.
\end{itemize}

Note that Stage I lasts at most $|X|$
rounds, and in the meantime Dominator cannot claim more than $a|X|$ vertices. Hence, for 
Stage II of Staller's strategy we have $|X_A'|\geq 4a|\calV(\calH)|$ for every $A\in \calV(\calH)$.
Stage II lasts at most $|\calV(\calH)|$ rounds,
since Staller only wants to claim one vertex in each of the sets $X_A'$. Hence, during this whole stage Dominator does not claim more than $a|\calV(\calH)|$ vertices, which implies that at any point in Stage II, there are at least $2a|\calV(\calH)|$
free vertices in each set $X_A'$ and thus
Staller can always claim a free vertex as required by Stage II of the strategy.  


It remains to prove that, by following the strategy, Staller 
prevents Dominator from getting a dominating set of $G$ 
of size at most $s$ within
the first $t$ rounds.
For contradiction, assume the opposite, i.e. 
Dominator claims a dominating set $D$ with $|D|\leq s$ within at most $t$ rounds.
Then $D' := D\cap X$ would be a dominating set by the following reason:
Because of Stage II, for each $A\in \mathcal{V}(\calH)$,
there is vertex $v\in X_A$ such that $v\notin D$.
Hence, since $N_G(v)=A\subset X$, 
we get $D\cap A \neq \varnothing$
and therefore $D'\cap A \neq \varnothing$. 
Thus, $D'$ dominates every vertex in $X_A$ for every $A\in \mathcal{V}(\calH)$,
and it also dominates every vertex in $X\setminus D' = X\setminus D$
since $G[X]$ is a complete graph. \\
By Claim~\ref{clm:transference.claim3}
it follows that there is some $F\in \calF$ with 
$F\subset D'\subset D$,
and note that $|F|\leq |D|\leq s$ by assumption. 
However, by following Stage~I, Staller ensures that Dominator needs to play at least $t$ rounds
on $X$ to occupy a winning set $F\in \calF$
of size at most $s$, a contradiction.

\medskip


\textbf{(W)} "$\Rightarrow$" The proof is essentially the same as for (M).

"$\Leftarrow$" The proof is very similar to the corresponding proof for (M).
Assume that Client has a strategy $\mathcal{S}_C$ in the 
Waiter-Client game on $\calH$ which prevents Waiter from 
getting a winning set of size at most $s$ 
within the first $t$ rounds. 
Now, in the Waiter-Client domination game on $G$,
Staller can apply the following strategy which distinguishes
three cases.
\begin{itemize}
\item[] \textbf{Case I:} When Dominator offers two elements of $X$,
then Staller distributes these vertices according to the strategy $\mathcal{S}_C$.
\item[] \textbf{Case II:} When Dominator offers a vertex $v\in X$ and a vertex $w \notin X$, then Staller claims $v$ and gives $w$ to 
Dominator. (For the strategy $\mathcal{S}_C$ in Case I, Client imagines that $v$ was not played yet.)
\item[] \textbf{Case III:} When Dominator offers two vertices 
$v\in X_{A_1}$ and $w\in X_{A_2}$ with $A_1,A_2\in \calV(\calH)$,
then Staller claims $v$ if he does not have a vertex in $X_{A_1}$ yet, and otherwise Staller claims $w$. The remaining vertex goes to Dominator.
\end{itemize}

By the strategy $\mathcal{S}_C$ in Case I,
and since Dominator never gets an element of $X$ in the other cases,
Dominator needs more than $t$ rounds in order to claim
a winning set of $\mathcal{F}$ with size at most $s$.
Moreover, by the strategy from Cases~II and~III,
Staller makes sure to receive at least one vertex in each set $X_A$.
For this note that for each $A\in \calV(\calH)$ there can be at 
most $|X|$ rounds of Case II and at
most $|\calV(\calH)|$ rounds of Case III, in which a vertex $x\in X_A$ is offered while Staller still does not have a vertex in $X_A$, but Staller decides to claim the other vertex which was offered by Dominator.

Now, using these observation, the same discussion as at the end of the proof of (M) gives that
Staller prevents Dominator from occupying a dominating set of $G$ of size at most $s$ within the first $t$ rounds.
\end{proof}

\end{document}
